%
% $Log: introduction.tex,v $
% Revision 1.4  2002/10/21 15:26:44  mjm
% to ohio, revised
%
% Revision 1.3  2001/08/15 18:07:03  mjm
% APPM preprint 466
%
% Revision 1.2  1999/08/09 19:55:27  mjm
% beta release
%
% Revision 1.1  1999/06/10 23:31:49  mjm
% Initial revision
% 
%

%\begin{probdescript}{Writing an Introduction to Your Problem}
\handout{Introductions and Conclusions}

\subsection*{Introductions}

The purpose of  an introduction is to explain the important
techniques used to solve your problem. The length may vary from a few
sentences to a paragraph. The length usually depends upon the number
of steps in your problem needed to get to a solution. A well
constructed introduction will make your answer and work easy to
follow. On the other hand, a poorly constructed introduction will
confuse the reader.  


We will illustrate some things to cover in your introduction with the
problem:
{\em
Use the Sandwich Theorem to find the
      asymptotes of the curve $y= \sin(x)/x$.} 


\begin{itemize}
\item  What the problem is about.
 Ask yourself: ``What am I doing?''
  Someone reading your problem should get an idea of what you are
  about to do. If there are multiple parts then explain each part.

\begin{badexample}
I'm going to calculate the asymptotes.
\end{badexample}

\begin{goodexample}
We will find
  the asymptotes of the function $y=\sin(x)/x$.\\
 \fbox{\rm This one supplies more information to the reader.} 
\end{goodexample}
  
\item The technique(s) you are going to use.
 This will usually correspond to
  the topic you are learning.  Your assigned problem will generally
  come from a section you are covering in class, and you can use this
as a guide.   Sometimes your
  problem will involve several concepts and you should mention what
  the predominant technique is that you are going to use. Do not
  mention the details of your calculations.

\begin{badexample}  
Calculating the limit of $y=\sin(x)/x$ will tell me where
  the asymptotes are.  
\end{badexample}

\begin{goodexample}
To find the asymptotes, we will take the limits as $x$ approaches
infinity or negative infinity. To determine these limits, we will use
the Sandwich Theorem.\\
\fbox{\rm The techniques employed are mentioned}
\end{goodexample}

\item Physical interpretation, if appropriate.
  Whether this
  is necessary depends upon the question. If applicable, this step is
  both helpful for you and the person who is reading your problem.
  Again, think about what you are doing. Give a physical description
  of the mathematical steps you are employing to solve your problem.
  If you were to draw a tangent line at a point, what would the
  picture look like?

\begin{badexample}  
By looking at the graph, $y$ goes to the $x$ axis.
\fbox{\rm What are you talking about?}
\end{badexample}

\begin{goodexample}
By looking at the graph of the function, we see that as $x$ approaches
$\pm\infty$, $y$ approaches~0, so the graph approaches the $x$ axis.\\
\fbox{\rm This description is much clearer.}
\end{goodexample}
 

\end{itemize}
Putting it
  all together: 
%\end{probdescript}

%\clearpage
%\secondpage
%\begin{probexamples}


      \begin{badexample}
        I'm going to calculate the asymptotes. Calculating the limit
        of $y=\sin(x)/x$ will tell me where the asymptotes are. By
        looking at the graph, y goes to the $x$ axis.
      \end{badexample}

      \begin{goodexample} We will find the asymptotes of the function
      $y=\sin(x)/x$.  To find the asymptotes, we will take the limits
      as $x$ approaches infinity or negative infinity. To determine
      these limits, we will use the Sandwich Theorem.  By looking at
      the graph of the function, we see that as $x$ approaches
      $\pm\infty$, $y$ approaches~0, so the graph approaches the $x$
      axis.  \end{goodexample}

\noindent\fbox{\parbox{\textwidth}{\rm 
        The bad introduction doesn't flow very well.  Someone
        reading it would not understand what you did.  It is obvious
        that you calculated the asymptotes, but when solving a problem,
        the techniques employed are as important as your answer. The
        good introduction briefly explains how you used the Sandwich
        Theorem along with the result you obtained. This would allow the
        reader to have a clear understanding of the mathematical
        calculations that follow.
}}

%\secondpage

%\subsection*{Examples:}

%{\bf Problem:}
%	{\em  Find $dy/dx$ if 
%      \begin{displaymath}
%        y=\int_1^{x^2} \cos(t)dt.
%      \end{displaymath}
%	}

%      \begin{badexample}
%        I'm going to take the derivative of this integral to find
%        $dy/dx$. Using the Chain Rule I found $dy/dx = 2x \cos(x^2).$
%    \end{badexample}

%\bigskip

%    \begin{goodexample}
%       I'm going to use the Fundamental Theorem of Calculus to
%      find $dy/dx$. By treating $y$ as a composite of 
%	$y=\int_1^u \cos(t)d$  and
%      $u=x^2$, I can use the Chain Rule to find $dy/dx = dy/du \times
%      du/dx$.  Evaluating these results in the answer $dy/dx = 2x
%      \cos(x^2)$.
%    \end{goodexample}

%\bigskip

%\noindent\fbox{\parbox{\textwidth}{\rm 
%      The bad example would not give the reader a clear guide to
%      follow your calculations.  In the good example, the
%      Fundamental theorem of Calculus is mentioned along with a
%      comment on how $y$ is going to be treated.  Thus when you do your
%      calculations, each step through the Chain Rule will have a
%      logical flow.
%}}


%\bigskip
%{\bf Problem:}
%	{\em 
%      The cost function at
%      XYZ Co. is $c(x)=x^3-6x^2+15x$. Is there a production level that
%      minimizes average cost? If so, what is it?  
%}

%\bigskip

%      \begin{badexample}
%        I'm going to calculate the derivative of the function to
%        find the marginal cost function. When this function equals the
%        average cost function we have a production level which
%        minimizes average cost.
%    \end{badexample}

%\bigskip

%    \begin{goodexample}
%      The marginal cost function is calculated as the first derivative
%      of the cost function.  The average cost function is $c(x)/x$. By
%      setting these two equations equal to each other and solving for
%      $x$, I get the production levels that might minimize the average
%      cost function. By performing the first and second derivative
%      tests on the average cost function at $x$ I can verify if the
%      production level $x$ is in fact a minimum.
%\marginpar{funny answer}
%    \end{goodexample}      

%\bigskip

%\noindent\fbox{\parbox{\textwidth}{\rm 
%        Clearly, this introduction provides a better outline about the
%        techniques used to solve the problem.
%}}


%\marginpar{conclusion}

\secondpage
\subsection*{Conclusions}

Your solution should end with a brief concluding paragraph. The
conclusion is a short, concise paragraph that summarizes the main
ideas, techniques and results discussed in the problem.  The
concluding paragraph should focus on the main results and main ideas
so that the reader can absorb the important parts of your solution to
the problem. The concluding paragraph should not
introduce any new material but simply bring back the reader to the main ideas
of your solution to the problem. This means that you should not
introduce something in the conclusion that was not discussed earlier in
your solution.

When you write a conclusion, put yourself in the potential reader's
position and imagine what kind of information this reader may look for
when he or she reads the concluding paragraph.  Also, ask yourself
``What do I want my reader to remember?''  Here are a few suggestions
of what you may want to include in the conclusion paragraph:

\begin{itemize}
\item Restate the problem you solved.
\item Restate the results and interpretation of the results. If your
  results are lengthy, such as a big table with numerical values,
  just  describe the results qualitatively. For example,
  give the range of values for your results.
\item Indicate if your results seem reasonable. If not, then explain why.
\item If you were unsuccessful in solving your problem, mention
  possible reasons for this.
\item If applicable, mention other problems that are related to the
  problem you solved.
\item Give suggestions for improvements, i.e., what else you could do
  for the problem.
\end{itemize}

These are just suggestions, and not all of them may be applicable for
your problem. 


%\secondpage

%\subsection*{Examples:}

{\bf Example Problem:}
  \emph{Find
        the dimensions of a right circular cylinder of volume
        $10^{-6}\ m^3$ such that a minimum amount of material is
        used.}  

\bigskip

      \begin{badexample}
        In this problem I used calculus to solve a mathematical
        problem.
\fbox{\rm This gives no information.}
    \end{badexample}

\bigskip

    \begin{badexample}
      I solved for $h$ in $\pi r^2h=10^{-6}$ and plugged the
      expression for $h$ into $2\pi r^2+2\pi rh$. Then I took the
      derivative of $2\pi r^2+\frac{2\times 10^{-6}}{r}$. Then I set
      the derivative equal to zero. Then I solved for $r$.
    \end{badexample}


\noindent\fbox{\parbox{\textwidth}{\rm 
      The last example fails to focus on the main ideas of the
      problem. It gives a (bad) summary of some of the steps used to
      solve the problem. It does not summarize the problem given to us
      and it does not clearly indicate what your bottom line was or
      if your were even successful in solving the problem.
}}

\bigskip

    \begin{goodexample}
      The problem discussed in this paper is to find the dimensions of
      a right circular cylinder of volume $10^{-6}\ m^3$ such that a
      minimum amount of material is used. The problem was solved by
      minimizing the area of the cylinder using the first and second
      derivative test. The optimal dimensions were found to be radius
      $\cong 5.42\ cm$ and height $\cong 10.84\ cm$. The method used
      for solving the problem can be applied to other optimization
      problems, e.g., finding the dimensions of other shapes with a
      fixed volume such that the used material is minimized. A
      possible development of the technique used could be to consider
      optimization of a function of more than two parameters and more
      than one constraint.
    \end{goodexample}

\noindent\fbox{\parbox{\textwidth}{\rm 
      Note that the suggested conclusion requires that we actually did
      discuss possible applications and developments of the technique
      earlier in the report, something that may not be required for
      you in your course.
}}

%\marginpar{conclusion}